% ch7_conclusion.tex

\chapter{Conclusion et perspectives}
\label{ch:conclusion}

Ce rapport a présenté le portage de StreamKM++ depuis l'implémentation Flink proposée par Bitsakis (2018)~\cite{bitsakis2018} vers Apache Spark. Les principaux composants algorithmiques, notamment l'algorithme de Lloyd pondéré, l'initialisation k-means++, la construction du coreset et la stratégie \emph{merge-and-reduce}, ont été implémentés indépendamment de Spark, puis intégrés à une couche distribuée reposant sur les RDD, \texttt{treeReduce} et Structured Streaming.

L'architecture retenue sépare les composants algorithmiques des mécanismes propres à Spark. Dans la version distribuée, chaque partition construit un coreset local à partir des observations qui lui sont attribuées. Ces résumés sont ensuite fusionnés afin de produire le coreset utilisé pour déterminer les centres finaux. Structured Streaming reprend ce principe pour traiter des micro-batches successifs.

L'évaluation réalisée sur les 581,012 observations de Covertype montre que la version distribuée produit des résultats très proches de ceux de l'implémentation séquentielle. Pour les valeurs de \(k\) étudiées, l'écart relatif entre leurs SSE reste inférieur à \(0.7\,\%\). Lors de la variation de la taille du coreset, les différences restent également faibles, avec un écart maximal d'environ \(0.63\,\%\).

La comparaison avec Spark MLlib montre par ailleurs que les SSE obtenues avec StreamKM++ restent proches de celles d'un K-Means appliqué directement aux données distribuées. Ces résultats indiquent que, dans les configurations testées, la construction distribuée du coreset avec Spark conserve une qualité de clustering proche de celle obtenue avec l'implémentation séquentielle.

Les conclusions doivent toutefois être replacées dans les limites du protocole expérimental. Chaque configuration a été exécutée une seule fois avec une graine fixe. Les expériences permettent donc de comparer les méthodes dans des conditions identiques, mais elles ne permettent pas de mesurer précisément la variabilité liée aux étapes aléatoires de l'algorithme.

Une expérimentation complémentaire a également porté sur la représentation mémoire des points. La variante SoA fondée sur un \texttt{DoubleBuffer} hors tas s'est révélée moins performante que la représentation AoS initiale pour les opérations évaluées. La représentation initiale a donc été conservée.

\section{Perspectives}

Une première perspective consiste à compléter l'évaluation de la qualité en répétant les expériences avec plusieurs graines. Il serait alors possible de calculer la moyenne et la dispersion des SSE obtenues et de mieux évaluer la stabilité des résultats.

L'influence de la taille du coreset pourrait également être étudiée avec davantage de valeurs de \(m\). Les expériences réalisées montrent une amélioration de la SSE lorsque la taille du coreset augmente, mais cette amélioration devient progressivement plus faible. Des essais supplémentaires permettraient de mieux caractériser le compromis entre la taille du résumé et la qualité du clustering.

Une étude spécifique des performances de l'implémentation distribuée constituerait également un prolongement important. Le projet s'est principalement concentré sur la qualité des résultats et n'a pas évalué systématiquement l'évolution du temps d'exécution en fonction du volume de données, du nombre de partitions ou des ressources disponibles.

Enfin, les travaux sur la représentation mémoire pourraient être poursuivis en évaluant une variante SoA utilisant un tableau \texttt{Array[Double]} plat conservé sur le tas. Cette expérience permettrait de distinguer plus clairement l'effet de l'organisation des coordonnées de celui de l'utilisation d'un buffer hors tas.

Dans son périmètre actuel, ce projet aboutit ainsi à une implémentation de StreamKM++ adaptée à Apache Spark et évaluée sur un jeu de données réel. Les résultats obtenus montrent que la construction distribuée du coreset conserve une qualité de clustering proche de celle de l'implémentation séquentielle dans les configurations étudiées.
